\chapter{ML-DSA Key Provisioning for PoC}

\section*{Purpose}
This chapter documents the temporary key-material strategy used for the current ML-DSA-OSH attachment phase.

\section{Current PoC Strategy}
The current proof of concept uses static secret-key material for controlled bring-up only. The project-owned ML-DSA-OSH shim loads a preformatted ML-DSA-87 secret key image from a local memory file and presents the expected upstream key segments to the imported core interface.

\section{Key Source}
The current PoC secret-key image is derived from the first ML-DSA-87 signing key vector in the imported upstream `KAT/SigGen\_sk\_87.txt` dataset. The project stores the derived byte-oriented memory image in `hw/include/mldsa\_osh\_poc\_sk\_87.mem`.

\section{Why This Is Acceptable for the Current Phase}
The current phase is focused on integrating the real signing core behind the stable engine adapter without changing the wrapper-visible contract. Static key material reduces bring-up variables and allows the adapter and shim path to be reviewed independently from production key-management design.

\section{Current Limitations}
\begin{itemize}[leftmargin=*]
  \item The static key image is suitable only for controlled development and simulation-oriented bring-up.
  \item The design does not yet provide secure key loading, lifecycle management, or hardware-backed secrecy guarantees.
  \item The current approach is intentionally incompatible with production security expectations.
\end{itemize}

\section{Replacement Direction}
A later phase shall replace the static key image with a real key-provisioning path. Candidate future approaches include PS-assisted loading into protected PL storage, device-bound provisioning, or other architecture-reviewed mechanisms consistent with the target platform and threat model.

\section{Documentation Requirement}
Any change to the key source, format, or provisioning path shall update this chapter, the hardware microarchitecture document, the ML-DSA-OSH integration guide, and the verification plan in the same workstream.